% Generated draft: requires mathematical and source review.
\documentclass[11pt]{article}
\usepackage[margin=25mm]{geometry}
\usepackage{fontspec}
\setmainfont{Latin Modern Roman}
\usepackage{amsmath,amssymb,mathtools,unicode-math}
\setmathfont{Latin Modern Math}
\usepackage{enumitem,needspace,xurl,hyperref}
\setlength{\parindent}{0pt}
\setlength{\parskip}{5pt}
\setlength{\emergencystretch}{3em}
\newcommand{\N}{\mathbb{N}}
\newcommand{\Z}{\mathbb{Z}}
\newcommand{\Q}{\mathbb{Q}}
\newcommand{\R}{\mathbb{R}}
\newcommand{\C}{\mathbb{C}}
\newcommand{\F}{\mathbb{F}}
\newcommand{\A}{\mathbb{A}}
\newcommand{\PP}{\mathbb{P}}
\newcommand{\E}{\mathbb{E}}
\newcommand{\eps}{\varepsilon}
\newcommand{\dd}{\,\mathrm{d}}
\begin{document}
\section*{UnsolvedMath ID: 20002087}
\textbf{Source problem code:} catalogue entry
\subsection{Definitions and mathematical statement}
Question \(1.29\), written in 2003, asks for the conjectured equivalence of Ozsv\'ath--Szab\'o theory and Seiberg--Witten theory, and notes Yi-Jen Lee's proposed approach. The word equivalence can mean groups, modules, exact systems, chain constructions or a cobordism-natural TQFT; the source does not fix one of these strengths.

The aid reports that Kutluhan--Lee--Taubes established an integral, relatively graded, module-compatible equivalence of the three-flavor Heegaard Floer exact system with the corresponding balanced monopole Floer system for every closed oriented three-manifold \(Y\) and fixed \(\operatorname{Spin}^c\) structure \(\mathfrak s\). The Heegaard flavors are \(HF^-,HF^\infty,HF^+\); the corresponding monopole flavors require the standard matching of orientation, duality and notation. The stated compatibility includes the coefficient/module actions, conventionally encoded by \(\Z[U]\otimes\bigwedge^*(H_1(Y;\Z)/\mathrm{torsion})\), and the exact-system maps. Precise flavor and grading conventions must be checked in the primary theorem. This resolves the source's group-system interpretation, not automatically every stronger functorial interpretation.

The aid says that a full cobordism-map/TQFT naturality theorem is not supplied merely by that equivalence statement and was not verified in the primary literature checked. This is a qualified scope assessment, not a claim of a new universal impossibility or an independently updated open status.

Its formal comparison criterion is: after fixing coefficients, orientation, \(\operatorname{Spin}^c\), flavor, duality and grading conventions, let
\(\Phi_Y,\Psi_Y:HF(Y)\to HM(Y)\)
be two compatible equivalence families and put
\[
 \Delta_Y=\Psi_Y^{-1}\Phi_Y\in\operatorname{Aut}(HF(Y)).
\]
For a fixed decorated cobordism \(W:Y_0\to Y_1\), write \(F_W\) and \(G_W\) for the matched HF and HM maps. If \(\Phi_{Y_1}F_W=G_W\Phi_{Y_0}\), then
\[
 \Psi_{Y_1}F_W=G_W\Psi_{Y_0}
 \quad\Longleftrightarrow\quad
 \Delta_{Y_1}F_W=F_W\Delta_{Y_0}.
\]
Thus the residual question is an explicit automorphism-commutation condition, not a consequence of isomorphic groups alone.

For the mod-two plus tower at \(S^3\), with its standard absolute grading,
\[
 HF^+(S^3;\F_2)\cong\mathcal T^+
       =\F_2[U,U^{-1}]/U\F_2[U],\qquad\deg U=-2,
\]
every grading-degree-zero \(\F_2[U]\)-module automorphism is the identity. Hence a discrepancy satisfying exactly those conventions is forced to be trivial there. This single-object normalization does not prove the commutation condition for all cobordisms or manifolds.

Finally a \(U\)-equivariant module isomorphism restricts to isomorphisms of \(\ker U^m\) and of their successive layers
\[
 \ker U^m/\ker U^{m-1}\qquad(m\geq1).
\]
The aid's transfer statement is in this module setting; it does not identify these layers with an additional Floer flavor without further comparison data. The appended ``Fifth day'' is the next heading, not part of the equivalence question.
\subsection{Short English statement}
The aid reports the closed-three-manifold group/module/exact-system HF--HM equivalence, while distinguishing full cobordism naturality. It reduces comparison of two equivalence families to a discrepancy commuting with cobordism maps. A mod-two graded plus-tower discrepancy is trivial at \(S^3\), and \(U\)-kernel layers transfer under a module isomorphism.
\subsection{Sources}
\begin{itemize}
\item Local UnsolvedMath registry, permanent ID \(20002087\), source code \texttt{catalogue entry}; canonical collection \texttt{aim-geometry-notes.json}, notes[424].
\item AIM workshop list ``Holomorphic curves in contact geometry,'' version 15 October 2003, Question \(1.29\): \url{https://aimath.org/WWN/contactgeom2/contactgeom2.pdf}. The source names Ozsvath--Szabo theory, Seiberg--Witten theory and Yi-Jen Lee; ``Fifth day'' is neighboring material. No PDF was opened in this session.
\item Kutluhan--Lee--Taubes are credited by the aid for the integral relatively graded three-flavor equivalence. Full authoritative titles, theorem numbers and exact notation conventions are not supplied. Original remarks and literature fields are empty, with no proposer or year.
\item Supplied research attempt \(1\), \texttt{partial\_result}, run status \texttt{partially\_solved}, rigor label ``proved,'' novelty confidence ``low.'' Registry classification \texttt{PARTIAL-PROGRESS}, updated 14 August 2026. The 2003 ``big open problem'' wording is historical, and the stronger naturality assessment is limited to the supplied research check.
\end{itemize}
\end{document}
